Three-dimensional object detection (3DOD) is the perception sub-task of producing, for every traffic agent in a scene, an oriented 3D bounding box together with a class label, and increasingly an instance velocity vector. It sits at the centre of the autonomous-driving stack: tracking, motion prediction, behaviour planning, and trajectory optimisation are all conditioned on the cuboids that 3DOD emits at typically 10 Hz. Compared with 2D object detection in image space --- where ground-truth boxes are axis-aligned rectangles in pixel coordinates --- 3DOD operates in metric Euclidean space, and the regression target is a seven-degree-of-freedom (7-DoF) tuple \texttt{(x,\ y,\ z,\ l,\ w,\ h,\ yaw)} in the ego-vehicle frame, sometimes extended to nine degrees of freedom by appending the planar velocity components \texttt{(\$v\_x\$,\ \$v\_y\$)} as required by the nuScenes Detection Score (NDS) protocol of Caesar et al.~(2020). Because actuators drive metric distances and brake decisions hinge on absolute time-to-collision, the consequences of localisation error in 3DOD are categorically more severe than in 2D detection. Formally, given a synchronised set of sensor observations at frame \texttt{t} consisting of one or more LiDAR sweeps \texttt{\$P\_t\$\ =\ \{(\$x\_i\$,\ \$y\_i\$,\ \$z\_i\$,\ \$r\_i\$)\}} of order \texttt{10\^{}5} points, a multi-camera RGB tuple \texttt{\$I\_t\$\ =\ \{\$I\_t\^{}k\$\}\_\{k=1..K\}} with known intrinsics \texttt{\$K\_k\$} and extrinsics \texttt{{[}\$R\_k\$\textbar{}\$t\_k\${]}} relative to the ego frame, and optionally a radar tensor \texttt{\$R\_t\$} that contains range-Doppler bins, a 3D o\ldots{}
